\documentclass[10pt]{article}
\usepackage[margin=1.5cm]{geometry}
\usepackage[T1]{fontenc}
\usepackage{lmodern}
\usepackage{microtype}
\usepackage{booktabs}
\usepackage{enumitem}
\usepackage{titlesec}
\usepackage[hidelinks]{hyperref}

\titlespacing*{\section}{0pt}{5pt}{1pt}
\titleformat*{\section}{\normalsize\bfseries}
\setlist{nosep,leftmargin=1.2em}
\setlength{\parskip}{2pt}
\setlength{\parindent}{0pt}
\setlength{\textfloatsep}{6pt}

\title{\vspace{-1.6cm}\large\textbf{Retrieval or Selection? Diagnosing and Improving XBRL Concept Linking\\
for the Open Financial LLM Leaderboard}}
\author{\normalsize Dev Gaglani \quad Stavan Desai}
\date{\normalsize Research proposal, November 2026}

\begin{document}
\maketitle
\vspace{-0.9cm}

\section{Motivation and gap}
Every number in an SEC filing is tagged in XBRL with one of about 17{,}400 US-GAAP 2024 concepts. Posed as direct
classification over this label set, large models fail: FinTagging reports zero F1 for GPT-4o and
DeepSeek-V3~[1]. The workable formulation is \emph{retrieve, then select}: a search step proposes candidate concepts
and an LLM chooses one. FinTagging reports that retrieval finds the correct concept among 200 candidates only about
33\% of the time, yet its own Table~8 shows that handing the model the gold candidates adds only about two points
(GPT-4o 0.183 to 0.202; Llama-3.1-8B reaches 0.108). It is therefore unclear which stage limits accuracy. Our first runs
suggest both do: on 484 FinCL items, when the gold concept is among the candidates, models select it only 32--44\% of
the time, though those items are the easier ones. Separately, public XBRL
datasets overlap and sometimes disagree on gold answers~[1--3], and the Open Financial LLM Leaderboard (OFLL)~[5] has no
XBRL or SEC-filing tasks.

\section{Research question}
\textbf{RQ:} \emph{How much of XBRL concept-linking error comes from candidate retrieval and how much from candidate
selection, and how much can cheap taxonomy constraints and a small fine-tuned dense retriever recover without oracle
information, for concepts both seen and unseen in training?}
\begin{itemize}
\item \textbf{H1 (diagnosis):} we measure coverage (gold among candidates) and selection given coverage per model. We
predict that even with the gold concept inserted at a random position, every model selects it less than 50\% of the time,
so selection limits accuracy as much as retrieval.
\item \textbf{H2 (retrieval):} without the given item type, the fine-tuned retriever with a taxonomy-type constraint raises
coverage at 100 over BM25 restricted to numeric types (31.2\%) and over an untrained dense retriever, on unseen as well as
seen concepts.
\item \textbf{H3 (transfer):} against that same no-oracle BM25 baseline, the paired 95\% CI of the accuracy difference
excludes zero for at least three of the four models; given paired half-widths near 2 points, this needs coverage at
100 to rise by about 7 points, our success threshold.
\end{itemize}
\textbf{If H2 holds but H3 fails,} selection is the binding constraint; that is itself an answer to the RQ, and the
report would then focus on the selection stage.

\section{Scope, data and overlapping questions}
\textbf{Main task:} FinCL concept linking from FinTagging's public release (491 test items, 484 scored; 93\% are table
rows). \textbf{Reported separately:} the COLING-25 tag queries (90 items; company and concept name to tag), which are a
different task. Fact extraction (FinNI, 199 test items) is used only to predict item types. Financial calculation and
filing QA are outside v1. The data come from four pinned public sources~[1--4], each with its SHA-256 hash and licence.
\begin{itemize}
\item \textbf{Overlaps:} records join only on a shared source id or identical normalised text (944 text and 248 id
joins; 11{,}401 records become 10{,}457 unique items across all sources; FinCL's 500 released items become 491).
MinHash near-duplicates~[7] go to human review, never automatic
merging (6 pairs). Each cluster of related items stays in one split.
\item \textbf{Conflicts:} 28 same-input, different-gold cases are flagged; same-dataset conflicts are excluded from the
primary score (7 FinCL items, leaving 484), cross-source ones keep the publisher's gold.
\item \textbf{Training data and leakage:} 24{,}798 fact--concept pairs are rebuilt from FinNI training documents (not an
official FinCL training set); 15 training items that conflict with or nearly duplicate test items are dropped. For 374
of the 484 test items the gold concept appears in training, and 132 of the 451 table rows share a row label with a
training row (83 with the same concept). We therefore report seen-concept, unseen-concept (110 items, CIs about
$\pm$9 points) and seen-row-label strata, and select models on a development split carved from training. Nothing is
tuned on test items.
\end{itemize}

\section{Method}
\textbf{Retrieval variants,} each scored for coverage at 1, 50, 100 and 200: (a) BM25~[6]; (b) BM25 restricted to
numeric types (no given type); (c) BM25 filtered on the type predicted by fact extraction, falling back to (b) when the
fact is not extracted; as oracle conditions, BM25 filtered on the given type, and a BM25 analogue of the soft type match
in FinTagging's released text-search script (type terms added to query and concept text; an approximation, not a
replication); (d) an untrained open dense retriever (e.g.\ bge-small, MIT licence), with and without these
filters; (e) the same model fine-tuned contrastively on the training pairs, with same-type sibling concepts as hard
negatives. FinTagging's paper reports an API embedding model (text-embedding-3-small), for which (d) is a weaker, local
stand-in. \textbf{Selection:} the same four LLMs choose from the top 100 candidates with a fixed prompt, so differences
come from retrieval alone. \textbf{Gold-candidate runs} insert the gold concept at a random position, which measures
selection on all items, not only the easy covered ones, and matches FinTagging's rerank-only setting.
\textbf{Preliminary evidence} (BM25, coverage at 200): 31.2\% unfiltered, 35.5\% with numeric types only, 36.8\% with the
soft type match and 41.5\% with the hard filter on the given type. On the 51 non-monetary items the soft match reaches
65\% and the hard filter 88\% (from 35\%); the 433 monetary items stay near 36\% under the hard filter but drop to
33.5\% under the soft match, below numeric-only. Always guessing ``monetary'' (right for
89\% of items) reaches only 32.2\%, so a predicted type must get the non-monetary items right.

\section{Evaluation}
\textbf{Models:} Qwen2.5-3B and 7B, Llama-3.1-8B and Mistral-7B (instruction-tuned, 4-bit, deterministic decoding,
6\,GB GPU). \textbf{Metrics:} coverage at $k$, selection given coverage, exact concept accuracy, 95\% bootstrap CIs
(Wilson at zero), paired bootstrap on identical items, and monetary and non-monetary items reported separately. LLM
decoding is greedy, so the second seed applies to retriever fine-tuning. FinTagging's numbers (different items) are shown apart.
\begin{table}[h]
\centering\small
\begin{tabular}{lcccc}
\toprule
 & \multicolumn{3}{c}{FinCL (484 items)} & Tag queries (90) \\
\cmidrule(lr){2-4}\cmidrule(lr){5-5}
Model & Closed-book & BM25 + given type, top 100 & Selected if covered & BM25, top 100 \\
\midrule
Qwen2.5-7B   & 1.2 [0.4, 2.3] & 15.5 [12.2, 18.8] & 43.9 & 36.7 [27.8, 46.7] \\
Llama-3.1-8B & 1.7 [0.6, 2.9] & 15.5 [12.4, 19.0] & 43.9 & 31.1 [21.1, 41.1] \\
Qwen2.5-3B   & 1.0 [0.2, 2.1] & 12.2 [9.3, 15.3]  & 33.9 & 26.7 [17.8, 35.6] \\
Mistral-7B   & 0.0 [0.0, 0.8] & 11.4 [8.5, 14.3]  & 32.2 & 14.4 [7.8, 22.2] \\
\bottomrule
\end{tabular}
\caption{\footnotesize Preliminary accuracy (\%, seed 0, 95\% CIs). Retrieval is deterministic, so coverage is the same
for every model: 35.3\% on FinCL (given type, an oracle) and 54.4\% on tag queries. Mistral closed-book is 0 of 484
(Wilson interval). Selection on covered items is far above FinTagging's rerank-only results; the gold-insertion runs
will show how much of that gap comes from covered items being easier.}
\end{table}

\section{Contribution to the Open Financial LLM Leaderboard}
The OFLL backend reads one results dataset and hard-codes five Greek-task columns; rows are keyed on hub model id,
revision and precision, with no fields for evaluation tracks or confidence intervals. \textbf{Reuse:} its React interface
(table, filters, column selector), FastAPI backend, Hugging Face result storage and model metadata schema.
\textbf{Our leaderboard page:} we host an XBRL leaderboard built on this codebase, showing coverage, selection and
accuracy per retrieval variant with CIs. \textbf{First pull request:} replace the hard-coded columns in the backend and
frontend with a task-group configuration, including optional fields for track and CI, so any instance can declare its
own tasks. Follow-up pull requests add a converter from evaluation results to the row format and update the outdated
contribution guide. Outside pull requests have waited months for review, so we present the design at the weekly
meeting first, and follow the DCO and Apache-2.0 rules.
An upstream XBRL task group follows once maintainers agree and our rows use pinned Hugging Face model revisions.

\section{Deliverables, timeline and roles}
\textbf{The first web page} shows where each model loses accuracy (retrieval or selection) and how each retrieval
variant changes it. \textbf{The report} answers the RQ and releases the deduplicated benchmark (data card, seen and
unseen splits) with the evaluation code. \textbf{Risks:} the text-row subset is small (33 items), so we make no
table-versus-text claims; results on the public subset may not hold on FinTagging's gated full set; FinanceBench
(CC-BY-NC) stays out of anything hosted; FASB definitions are left out until their terms of use are checked.

\begin{table}[h]
\centering\small
\begin{tabular}{p{1.1cm}p{10.3cm}p{2.6cm}}
\toprule
Week & Work & Lead \\
\midrule
1    & Numeric-only baseline, gold-candidate and predicted-type runs; development and seen/unseen splits. Task-group design agreed at the
weekly meeting & Dev; Stavan \\
2--3 & Dense baseline, type filters and fine-tuned retriever. Configuration PR, then result converter & Dev; Stavan \\
4    & End-to-end runs for all variants and models; second fine-tuning seed. Our leaderboard page & Dev; Stavan \\
5--6 & Analysis and report; data card and code release & Both \\
\bottomrule
\end{tabular}
\end{table}
\textbf{Roles.} \emph{Dev Gaglani} leads the model runs, retriever training and experimental analysis. \emph{Stavan
Desai} leads the leaderboard contribution and page, the dataset documentation and the seen/unseen and development
splits. Both write and review the report.

\vspace{3pt}
{\footnotesize\textbf{References.}
[1] Y.~Wang et al. FinTagging: Benchmarking LLMs for extracting and structuring financial information. arXiv:2505.20650, 2025.
[2] K.~Wang et al. FinNLP-FNP-LLMFinLegal-2025 shared task: Regulations Challenge. COLING workshop, 2025.
[3] S.~Han et al. XBRL Agent: Leveraging LLMs for financial report analysis. ICAIF, 2024.
[4] P.~Islam et al. FinanceBench: A new benchmark for financial question answering. arXiv:2311.11944, 2023.
[5] FINOS and SecureFinAI Lab. Open Financial LLM Leaderboard. \url{https://github.com/finos-labs/Open-Financial-LLMs-Leaderboard}.
[6] S.~Robertson and H.~Zaragoza. The probabilistic relevance framework: BM25 and beyond. 2009.
[7] A.~Z.~Broder. On the resemblance and containment of documents. 1997.\par}
\end{document}
